\documentclass[11pt]{article}
\usepackage[T1]{fontenc}
\usepackage[utf8]{inputenc}
\usepackage{lmodern}
\usepackage[french]{babel}
\usepackage[a4paper,margin=2.6cm]{geometry}
\usepackage{amsmath,amssymb,amsthm}
\usepackage{tikz}
\usetikzlibrary{arrows.meta,positioning,fit,backgrounds}
\usepackage{booktabs}
\usepackage{microtype}
\usepackage[colorlinks=true,linkcolor=black,citecolor=black,urlcolor=blue!50!black]{hyperref}
\usepackage{xurl}

\newtheorem{proposition}{Proposition}
\newtheorem{definition}{Définition}
\newtheorem{corollaire}{Corollaire}

\title{Une porte d'admission à décroissance pondérée, à escompte
  d'indépendance des sources et à seuil adaptatif de charge, pour un
  étage d'argumentation NP-complet}
\author{Note d'ingénierie de prototype --- \texttt{moderation-pipeline-prototype}}
\date{Août 2026}

\begin{document}
\maketitle

\begin{abstract}
Un pipeline d'escalade qui bascule vers l'énumération d'extensions préférées
dans un cadre d'argumentation abstraite (Dung, 1995) lorsque la sémantique
fondée ne tranche pas hérite de la complexité, dans le pire cas, de ce
recours~: l'acceptation crédule d'une extension préférée est NP-complète.
Une règle d'admission qui ouvre l'accès à cet étage coûteux dès qu'un
dossier stagne et que le délai d'attente expire est exploitable au coût
attaquant $O(1)$~: une seule salve d'arguments plausibles et mutuellement
attaquants suffit à forcer l'admission. Nous spécifions, implémentons et
soumettons à une évaluation adversariale une porte de remplacement,
construite par composition de trois éléments~: un score de pression
d'escalade propre à chaque dossier, qui décroît géométriquement à chaque
pas de temps (transposition directe de la règle d'évaporation des
phéromones utilisée en optimisation par colonies de fourmis), un
renforcement plafonné par pas et ramené à un poids quasi nul lorsqu'il
provient d'une source déjà vue dans une fenêtre récente (codant une
exigence de corroboration provenant de sources réellement indépendantes,
et non simplement d'un délai écoulé), et un seuil d'admission qui
lui-même s'élève avec le taux d'occupation courant de l'étage coûteux à
l'échelle du système (un signal agrégé, dépendant de la densité, dans
l'esprit du \emph{quorum sensing}). Nous dérivons sous forme close le
comportement en régime permanent et le temps minimal d'admission de cette
porte à partir de la récurrence linéaire qui la gouverne, et vérifions
cette dérivation contre la simulation déployée~: une salve en un seul pas
de temps, de taille quelconque (1 à 5000), est bloquée sans condition~;
une inondation persistante issue d'une même cohorte est bloquée
structurellement, son score en régime permanent étant provablement
incapable d'atteindre le seuil~; une inondation persistante utilisant à
chaque pas une source fraîche et d'apparence non liée n'est pas bloquée
mais voit son admission retardée du coût attaquant $O(1)$ à $O(5)$ tours
soutenus et diversifiés en identité, conformément à la prédiction
analytique à un pas près. Une recherche dédiée d'antériorités, en seize
requêtes, n'a trouvé aucun précédent direct pour cette composition
précise à trois éléments gardant l'accès à un étage d'argumentation
NP-complet~; les transferts biologiques pris isolément, et le domaine
général du contrôle d'admission, sont en revanche bien documentés, et
cités comme tels plutôt que passés sous silence. Le résultat est
rapporté comme une défense partielle honnêtement bornée --- une
augmentation réelle et quantifiée du coût attaquant, non une clôture ---
cohérente avec le résidu identifié pour la variante la plus difficile, celle
d'un attaquant patient et diversifiant son identité.
\end{abstract}

\section{Le problème}
\label{sec:problem}

Un pipeline de modération de contenu à quatre étages achemine vers un
cadre d'argumentation abstraite (\emph{abstract argumentation framework},
AF) \cite{dung1995} les dossiers sur lesquels un étage de classification
sélective s'est abstenu~: les arguments et les attaques sont construits
à partir des déclarations et réfutations brutes du dossier, et
l'\emph{extension fondée} du cadre --- le plus petit point fixe de la
fonction caractéristique $F_{AF}$, atteignable en temps polynomial --- en
constitue la sortie principale, tractable. Lorsque deux groupes
d'arguments qui s'attaquent mutuellement laissent l'extension fondée vide
(le motif dit du « diamant de Nixon »), le pipeline se rabat sur
l'énumération des \emph{extensions préférées}, dont le problème
d'acceptation crédule est NP-complet \cite{fichte2024}. Ce recours est
une échappatoire nécessaire, délibérément bornée, non un défaut de
conception~: sans lui, un dossier bloqué ne se résoudrait jamais.

La vulnérabilité se situe à la frontière de l'\emph{admission}, non dans
la machinerie d'argumentation elle-même. Si un dossier stagnant est admis
à l'étage coûteux dès que son délai d'escalade expire, un attaquant
capable de produire, en une seule salve, des arguments plausibles et
mutuellement attaquants force l'admission au coût $O(1)$~: aucune
persistance, aucun signal d'identité distinguable, rien que du volume à
un instant donné. L'objectif de l'attaquant n'est pas de gagner le
débat --- la porte d'intégrité des preuves du scénario A ferme déjà
cette voie \cite{specA} --- mais de faire payer à la plateforme, à
répétition et à faible coût, la taxe NP-complète.

\section{La porte}
\label{sec:gate}

\subsection{Source de conception}

Deux mécanismes biologiques ont motivé la construction, utilisés ici
comme source génératrice de la conception, et non comme justification
apportée après coup à un mécanisme déjà bâti.

\textbf{Renforcement stigmergique.} Chez les fourmis en quête de
nourriture, les pistes de phéromones s'évaporent en continu~; la règle de
mise à jour canonique de Dorigo et Stützle s'écrit
$\tau \leftarrow (1-\rho)\tau$, appliquée à chaque itération
\cite{dorigo2004}. Les expériences du pont double de Deneubourg ont établi
le fait complémentaire que ce qui compte dans la force d'une piste est le
renforcement apporté par des fourmis explorant \emph{indépendamment} les
unes des autres, et non un seul individu repassant plusieurs fois sur son
propre chemin~; AntNet, de Di Caro et Dorigo, impose la même exigence
d'indépendance dans un cadre de systèmes distribués, en bâtissant sa
règle de renforcement sur une « communication indirecte et non
coordonnée » entre agents qui « collectent l'information en explorant le
réseau » \cite{dicaro1998}.

\textbf{Détection de quorum.} Chez les bactéries, le \emph{quorum
sensing} régule l'expression génique en fonction de la densité de
population, détectée par un signal agrégé, propre à l'environnement
partagé (la concentration d'un auto-inducteur), plutôt que par l'état
privé d'une cellule isolée \cite{miller2001}. L'idée structurelle
transposée ici est étroite et précise~: subordonner une admission
collective à un signal d'occupation partagé et agrégé, non au seul score
local d'un dossier pris isolément.

\subsection{Spécification formelle}

\begin{definition}[État d'admission]
Un dossier $c$ possède un score $S_c(t) \in \mathbb{R}_{\geq 0}$,
initialisé à $0$ au premier contact. À chaque pas de temps discret
$t = 0, 1, 2, \dots$, l'ensemble des événements de renforcement observés
pour $c$ est un multiensemble d'identifiants de cohorte-source
$R_c(t) \subseteq \mathcal{K}$.
\end{definition}

\begin{definition}[Contribution par pas]
Pour chaque cohorte $k \in R_c(t)$ (dédupliquée à l'intérieur du pas),
soit $\mathrm{frais}(k,c,t) = 1$ si $k$ n'a pas renforcé $c$ au cours des
$w$ derniers pas (la fenêtre d'indépendance), et $0$ sinon. La
contribution brute du pas est
\[
r_c(t) = \sum_{k \in R_c(t)} \big[ \mathrm{frais}(k,c,t)\cdot 1 +
\big(1-\mathrm{frais}(k,c,t)\big)\cdot \delta \big],
\qquad 0 < \delta \ll 1,
\]
plafonnée à un montant fixe par pas $\gamma$~: la contribution appliquée
est $\hat r_c(t) = \min(r_c(t), \gamma)$.
\end{definition}

\begin{definition}[Mise à jour du score]
\[
S_c(t) = (1-\rho)\, S_c(t-1) + \hat r_c(t), \qquad \rho \in (0,1].
\]
\end{definition}

\begin{definition}[Seuil d'admission adaptatif à la charge]
Soit $L(t)$ le nombre de dossiers occupant l'étage coûteux au pas $t$ (un
signal partagé et agrégé, propre au système entier, non une propriété du
seul dossier $c$), $\theta_0$ le seuil de base, $\kappa$ la sensibilité à
la charge, et $C$ une capacité de référence. Le seuil d'admission est
\[
\Theta(t) = \theta_0 \Big(1 + \kappa \frac{L(t)}{C}\Big).
\]
Le dossier $c$ est admis à l'étage coûteux au premier pas $t_{\ast}$ tel
que $S_c(t_{\ast}) \geq \Theta(t_{\ast})$.
\end{definition}

Cette spécification compose trois mécanismes indépendamment motivés~: une
décroissance géométrique (évaporation stigmergique), une contribution
plafonnée par pas et escomptée selon l'indépendance des sources
(renforcement stigmergique exigeant la diversité des sources), et un
seuil adaptatif à la charge (détection de quorum). Chaque élément, pris
isolément, recombine un principe déjà connu --- décroissance plus
plafond par pas est structurellement un limiteur de débit de type
\emph{leaky bucket}/EWMA, et un seuil qui s'élève avec la charge est
structurellement proche de \emph{Random Early Detection}
\cite{floyd1993}. L'objet de la composition est la contribution
spécifique, revendiquée comme telle~: faire garder ensemble, par ces
trois mécanismes, l'accès à un étage d'argumentation NP-complet, en
exigeant un sourçage indépendant --- et non le simple écoulement du
temps --- pour qu'un renforcement compte à plein poids.

\subsection{Comportement en forme close}

Considérons le pire cas pour un attaquant capable de livrer, à chaque
pas, la contribution appliquée maximale $c_{\ast}$ (soit $\gamma$, si une
source d'apparence fraîche est disponible à chaque pas, soit
$\delta\gamma$, si la même source se répète). En déroulant la récurrence
à partir de $S_c(0) = 0$~:
\[
S_c(n) = c_{\ast} \sum_{i=0}^{n-1} (1-\rho)^i = c_{\ast}\,\frac{1-(1-\rho)^n}{\rho}.
\]

\begin{proposition}[Plafond en régime permanent]
\label{prop:plafond}
$\displaystyle \lim_{n\to\infty} S_c(n) = \frac{c_{\ast}}{\rho}$. Si
$\Theta \geq c_{\ast}/\rho$ à la charge en vigueur au moment du contact,
aucune suite de renforcements plafonnés au taux $c_{\ast}$ ne peut jamais
admettre $c$, quelle que soit sa durée.
\end{proposition}
\begin{proof}[Argument]
Immédiat à partir de la forme close~: $(1-\rho)^n \to 0$ quand
$n\to\infty$ pour $\rho \in (0,1]$, si bien que $S_c(n)$ croît de façon
monotone vers $c_{\ast}/\rho$ par valeurs inférieures (chaque terme de la
somme est positif ou nul). Si le plafond n'atteint pas $\Theta$, aucun
$n$ fini ne le peut non plus, puisque $S_c(n) < c_{\ast}/\rho$ pour tout
$n$ fini.
\end{proof}

C'est ce mécanisme qui rend l'inondation persistante issue d'une même
cohorte \emph{structurellement} bloquée, et non pas simplement retardée~:
un renforcement répété depuis une seule source est escompté à
$c_{\ast} = \delta\gamma$, dont le plafond $\delta\gamma/\rho$ est conçu
pour rester sous $\theta_0$.

\begin{proposition}[Nombre minimal de pas jusqu'à l'admission]
\label{prop:minpas}
Si $\Theta < c_{\ast}/\rho$ (le plafond atteint bien le seuil), le
nombre minimal de pas jusqu'à l'admission, en supposant que l'attaquant
sature $c_{\ast}=\gamma$ à chaque pas depuis une source d'apparence
fraîche, est
\[
n_{\min} = \frac{\ln\!\big(1 - \Theta\rho/\gamma\big)}{\ln(1-\rho)}.
\]
\end{proposition}
\begin{proof}[Argument]
On résout $S_c(n) \geq \Theta$ pour le plus petit entier $n$, en
utilisant la forme close de la proposition~\ref{prop:plafond} avec
$c_{\ast}=\gamma$~: $\gamma(1-(1-\rho)^n)/\rho \geq \Theta \iff
(1-\rho)^n \leq 1-\Theta\rho/\gamma$. Comme $0<1-\rho<1$, prendre le
logarithme des deux membres inverse l'inégalité~:
$n\ln(1-\rho) \leq \ln(1-\Theta\rho/\gamma)$, et diviser par la quantité
négative $\ln(1-\rho)$ donne
$n \geq \ln(1-\Theta\rho/\gamma)/\ln(1-\rho)$.
\end{proof}

\begin{corollaire}[La tension d'étalonnage]
\label{cor:tension}
Notons $x = \Theta\rho/\gamma$ (évalué en $\Theta=\theta_0$, à charge
nulle). La condition d'admission d'un dossier authentique (une
probabilité moyenne de renforcement indépendant $p \gtrsim x$ par pas,
pour que le mécanisme reste utilisable en pratique) et la durée minimale
de l'inondation soutenue nécessaire à l'attaquant, $n_{\min}(x)$ selon la
proposition~\ref{prop:minpas}, sont gouvernées par la même quantité $x$.
Abaisser $x$ pour faciliter l'admission des dossiers authentiques
raccourcit nécessairement $n_{\min}$, et relever $x$ pour allonger
$n_{\min}$ durcit nécessairement le seuil pour les dossiers authentiques.
Aucun choix des constantes ajustables ne permet de relâcher l'un sans
resserrer l'autre.
\end{corollaire}

Le corollaire~\ref{cor:tension} est énoncé et présenté comme une
limitation authentique de cette classe de modèles, non résolue par le
réglage des paramètres~: c'est une propriété structurelle de toute porte
construite à partir d'une décroissance géométrique et d'un plafond fixe
par pas, indépendante des constantes particulières retenues.

\section{Évaluation adversariale}
\label{sec:eval}

La porte ($\rho=0{,}10$, $\gamma=1{,}0$, $w=3$ pas, $\delta=0{,}05$,
$\theta_0=4{,}0$, $\kappa=0{,}5$, $C=5$) a été confrontée à l'ancienne
référence statique --- admettre sans condition dès le premier
renforcement, une fois le dossier stagnant et son délai d'attente
expiré --- selon cinq scénarios adversariaux (code source et tests
complets dans le dépôt du prototype accompagnant ce document~; résultats
ci-dessous ré-exécutés indépendamment pour la rédaction du présent
article, \texttt{pytest tests/test\_admission\_control.py -v -s}, 6/6
réussis, 0{,}38\,s).

\begin{table}[h]
\centering
\small
\begin{tabular}{@{}p{4.0cm}p{4.6cm}p{4.6cm}@{}}
\toprule
Variante d'attaque & Référence statique & Porte défendue \\
\midrule
Salve en un seul pas, taille 1 à 5000, structure de cohorte quelconque
  & 100\,\% admis, pas 0 & \textbf{0\,\% admis, sans condition} \\
Inondation persistante, même cohorte, 200 pas
  & Admis au pas 0 & \textbf{Structurellement inatteignable} (plafond
    $0{,}5 < \theta{=}4{,}0$) \\
Inondation persistante, cohorte usurpée (source fraîche à chaque pas)
  & Admis au pas 0 & Admis au pas 4 (\textbf{retard $\times 5$}, conforme
    à $n_{\min}\!\approx\!4{,}85$) \\
Pulvérisation~: 200 dossiers, salve isolée chacun, 15 litiges
  authentiques concurrents
  & 215/215 admis (100\,\% de réussite pour l'attaque) &
    \textbf{15/215 admis} (0\,\% de réussite, 93{,}0\,\% d'admissions en
    moins) \\
12 campagnes soutenues synchronisées, adaptation à la charge activée ou
  non
  & --- & étalement sur 0 pas (désactivée) $\to$ \textbf{étalement sur
    13 pas} (activée) \\
\bottomrule
\end{tabular}
\caption{Résultats mesurés, référence statique contre porte défendue. Un
contrôle de cohérence (20 litiges authentiques, renforcés
organiquement) a confirmé l'admission de 20 dossiers sur 20 par la porte
défendue~: les chiffres de 0\,\% ci-dessus traduisent bien une
discrimination réelle, non un rejet en bloc.}
\end{table}

Le résultat obtenu contre la cohorte usurpée est le résidu honnête, non
vaincu~: un attaquant disposé à soutenir une campagne réellement
diversifiée en identité (des cohortes fraîches, jamais réutilisées, à
chaque pas) finit tout de même par réussir, à un coût relevé de $O(1)$ à
$O(n_{\min})$. Ce résultat est présenté comme une défense partielle, non
comme une clôture, exactement ce que prédit la
proposition~\ref{prop:plafond} pour tout taux de contribution $c_{\ast}$
dont le plafond dépasse $\Theta$.

Deux bogues d'implémentation, découverts en cours de test et signalés
plutôt que dissimulés, méritent d'être consignés comme des leçons
d'ingénierie transposables à cette classe de mécanismes~: (1) une
première version sommait les contributions d'une même cohorte
\emph{à l'intérieur} d'un seul pas avant déduplication, ce qui permettait
à une inondation d'une même cohorte d'atteindre le plafond par le seul
volume intra-pas, contournant l'escompte d'indépendance --- corrigé en
dédupliquant les cohortes par pas avant application du plafond~;
(2) une première version des tests adversariaux ne libérait jamais les
dossiers admis du signal de charge, produisant un seuil qui ne cessait de
croître et affamait les dossiers authentiques ultérieurs, pur artefact de
simulation --- corrigé en modélisant un temps de séjour borné à l'étage
coûteux, non en assouplissant l'assertion du test. La leçon générale~:
dans une porte de style \emph{quorum sensing}, le signal de charge doit
être lié à une occupation réelle et se vidant effectivement, non à un
compteur cumulatif d'admissions.

\section{Antériorités}
\label{sec:priorart}

Une recherche dédiée en seize requêtes (journal complet conservé dans le
document de spécification du dépôt accompagnant) a couvert le contrôle
d'admission d'inspiration biologique, la littérature sur le
\emph{quorum sensing} appliqué à l'informatique, la littérature sur le
déni de service contre les cadres d'argumentation et la complexité des
attaques ciblant l'extension préférée, ainsi que les systèmes
d'arbitrage et de justice décentralisée. Aucun article n'a été trouvé
composant les trois éléments ci-dessus (renforcement pondéré par
décroissance et escompté selon l'indépendance des sources~; plafond
strict par pas~; seuil adaptatif à la charge de type quorum) comme porte
d'admission à un étage d'argumentation NP-complet.

Les transferts pris isolément sont, en revanche, bien documentés, et
cités comme tels~: le renforcement par décroissance de phéromone relève
de l'optimisation par colonies de fourmis la plus classique
\cite{dorigo2004,dicaro1998}~; une combinaison décroissance-plus-plafond,
prise seule, est structurellement un limiteur de débit de type
\emph{leaky bucket}/EWMA, primitive standard et non biologique~; un seuil
d'admission qui s'élève avec la charge est structurellement proche de
\emph{Random Early Detection} \cite{floyd1993}~; le \emph{quorum sensing}
a déjà été transposé à l'informatique, pour le partitionnement en
grappes \cite{wang2013} et pour l'agrégation d'anomalies de sécurité
multi-agents \cite{fioravanti2012}, ni l'un ni l'autre n'étant une porte
d'admission pour un étage d'argumentation. Le résultat adjacent le plus
proche trouvé pour le patron général de transfert biologie-informatique
employé ici est une défense contre les inondations par système immunitaire
artificiel \cite{aiartnet2024}, qui repose sur la détection d'anomalies par
entropie du trafic et une mise en quarantaine immunitaire --- un
mécanisme différent, non fondé sur une décroissance de pression
d'admission.

\section{Menaces à la validité}
\label{sec:threats}

L'évaluation est une simulation synthétique à pas discrets, non un test
contre une dynamique d'arrivée réelle en temps continu~; la traduction de
la granularité du pas retenue ici vers l'unité de temps d'un déploiement
réel reste un travail d'ingénierie non vérifié. La condition de cohorte
usurpée accorde à l'attaquant une capacité généreuse (une identité
fraîche et structurellement indétectable à chaque pas)~; le plafond réel
de diversité d'identité qu'un détecteur effectif imposerait n'est pas
modélisé ici. Cette porte suppose que les couches de détection de
coordination et d'intégrité des preuves présentes ailleurs dans le
pipeline sont inefficaces dans le pire cas testé ici (la condition de
cohorte usurpée du scénario B)~; en pratique, ces couches relèveraient
encore le coût réel de l'attaquant, mais cette composition n'est pas
mesurée dans le présent document. Aucune affirmation n'est faite selon
laquelle cette porte, ou le pipeline dans lequel elle s'insère, résoudrait
le problème authentiquement ouvert de juger un contenu réellement inédit
au regard d'une spécification contestée --- une question orthogonale et
structurellement indécidable, que ce mécanisme ne touche pas.

\section{Conclusion}

La porte d'admission spécifiée ici ferme la variante la plus économique
et la plus réaliste d'une attaque par épuisement de ressources contre un
recours d'argumentation NP-complet --- des salves en un seul pas, de
taille quelconque, concentrées sur un dossier ou pulvérisées sur
plusieurs --- et relève, sans l'éliminer, le coût de la variante la plus
difficile, patiente et diversifiée en identité, de $O(1)$ à
$O(n_{\min})$ tours soutenus, sous la paramétrisation testée. Cette
composition n'a pas été retrouvée dans une recherche dédiée
d'antériorités à seize requêtes, tandis que ses transferts biologiques et
informatiques constitutifs ont chacun, individuellement, été trouvés et
cités comme précédents réels, non présentés comme novateurs. La tension
d'étalonnage du corollaire~\ref{cor:tension} est rapportée comme une
propriété inhérente à cette classe de modèles, non comme une limitation
propre aux constantes choisies ici.

\appendix
\section{Placement dans le pipeline}

\begin{center}
\resizebox{\textwidth}{!}{%
\begin{tikzpicture}[
  node distance=1.1cm and 1.6cm,
  box/.style={draw, rounded corners, minimum width=3.3cm, minimum height=1cm, align=center, font=\small},
  gate/.style={draw, rounded corners, thick, minimum width=3.6cm, minimum height=1.3cm, align=center, font=\small, fill=black!5},
  arr/.style={-{Stealth[length=2.2mm]}, thick}
]
  \node[box] (dung) {Étape 2b\\Résolveur de Dung\\(ext. fondée, temps $P$)};
  \node[box, right=of dung] (stale) {Ext. fondée vide\\+\\délai atteint};
  \node[gate, right=of stale] (gate) {Porte d'admission\\(cet article)\\décroiss. + plafond + quorum};
  \node[box, right=of gate, fill=red!5] (np) {Énum. ext.\\préférées\\(NP-complet)};

  \draw[arr] (dung) -- (stale);
  \draw[arr] (stale) -- (gate);
  \draw[arr] (gate) -- node[above,font=\footnotesize]{admis} (np);
  \draw[arr] (gate.south) to[out=250,in=290,looseness=1.6]
      node[below,font=\footnotesize] {rejeté : le score se redécroît, nouvel essai au pas suivant} (gate.south west);
\end{tikzpicture}%
}
\end{center}

\begin{thebibliography}{99}

\bibitem{dung1995} P.~M. Dung. On the acceptability of arguments and its
fundamental role in nonmonotonic reasoning, logic programming and
n-person games. \emph{Artificial Intelligence}, 77(2):321--357, 1995.

\bibitem{fichte2024} J.~K. Fichte, S.~A. Gaggl, D.~Rusovac, et al.
Complexity of claim-centric abstract argumentation frameworks.
\emph{AAAI}, 2024.

\bibitem{dorigo2004} M.~Dorigo et T.~Stützle. \emph{Ant Colony
Optimization}. MIT Press, 2004.

\bibitem{dicaro1998} G.~Di Caro et M.~Dorigo. AntNet: Distributed
stigmergetic control for communications networks. \emph{Journal of
Artificial Intelligence Research}, 9:317--365, 1998.

\bibitem{miller2001} M.~B. Miller et B.~L. Bassler. Quorum sensing in
bacteria. \emph{Annual Review of Microbiology}, 55:165--199, 2001.

\bibitem{floyd1993} S.~Floyd et V.~Jacobson. Random early detection
gateways for congestion avoidance. \emph{IEEE/ACM Transactions on
Networking}, 1(4):397--413, 1993.

\bibitem{wang2013} X.~Wang. A quorum sensing inspired algorithm for
dynamic clustering. \emph{arXiv:1303.3934}, 2013.

\bibitem{fioravanti2012} M.~E. Fioravanti~II. Digital quorum sensing for
self-organizing malware / a quorum sensing pattern for multi-agent
self-organizing security systems. Florida Institute of Technology,
2012--2013.

\bibitem{aiartnet2024} Adaptive artificial immune networks for
mitigating DoS flooding attacks. \emph{arXiv:2402.07714}.

\bibitem{specA} Spécification interne~: \emph{A four-stage
content-moderation escalation pipeline: formal specification, attack
analysis, and prototype verdict}, \texttt{moderation-pipeline-prototype/SPEC.md},
2026.

\end{thebibliography}

\end{document}
